\documentclass{article}
\usepackage{amsmath}
\usepackage{amsfonts}
\usepackage{amssymb}
\usepackage{geometry} % Required for customizing margins

\geometry{margin=1in} % Sets 1-inch margins on all sides


\title{EMCH 360 Assignment \#4}
\author{Jacob C. Vaught}
\date{October 22nd 2023}

\begin{document}

\maketitle



\section*{Problem \#1}
Determine the horizontal component of the force exerted by the blowing wind on the exhaust gases, given:
\begin{itemize}
    \item Diameter of the chimney \(D = 4 \, \text{ft}\)
    \item Speed of exhaust gases leaving the chimney \(V_1 = 6 \, \text{ft/s}\)
    \item Speed of the wind \(V_2 = 15 \, \text{ft/s}\)
    \item Density of air \(\rho_{\text{air}} = 2.373 \times 10^{-3} \, \text{slug/ft}^3\)
\end{itemize}

\subsection*{Assumptions}
\begin{enumerate}
    \item Steady flow
    \item One-dimensional flow
    \item No other forces acting on the fluid element
\end{enumerate}

\subsection*{Calculations}

\subsubsection*{Mass Flow Rate}
First, calculate the mass flow rate \(\dot{m}\):
\[
\dot{m} = \rho A V_1
\]
Area \(A\) of the chimney is:
\[
A = \pi \left( \frac{D}{2} \right)^2 = \pi \left( \frac{4 \, \text{ft}}{2} \right)^2 = 12.566 \, \text{ft}^2
\]
\[
\dot{m} = 2.34 \times 10^{-3} \, \text{slug/ft}^3 \times 12.566 \, \text{ft}^2 \times 6 \, \text{ft/s} = 0.176 \, \text{slug/s}
\]

\subsubsection*{Horizontal Component of Force}
Use the formula \( F_x = \dot{m} \times V_2 \):
\[
F_x = 0.176 \, \text{slug/s} \times 15 \, \text{ft/s} = 2.64 \, \text{lb}
\]
The horizontal component of the force that the blowing wind exerts on the exhaust gases is \(2.64 \, \text{lb}\).





\newpage
\section*{Problem \#2}
Determine the magnitude and direction of the anchoring force needed to hold the horizontal elbow and nozzle combination shown in Fig. P5.45 in place, given:
\begin{itemize}
    \item Diameter at Section (1), \(D_1 = 300 \, \text{mm}\)
    \item Velocity at Section (1), \(V_1 = 2 \, \text{m/s}\)
    \item Diameter at Section (2), \(D_2 = 160 \, \text{mm}\)
    \item Atmospheric pressure, \(P_{\text{atm}} = 100 \, \text{kPa}\)
    \item Gage pressure at Section (1), \(P_1 = 100 \, \text{kPa}\)
\end{itemize}

\subsection*{Assumptions}
\begin{enumerate}
    \item Steady flow
    \item One-dimensional flow
    \item Incompressible flow
    \item No losses due to friction or other factors
\end{enumerate}

\subsection*{Calculations}

\subsubsection*{Mass Flow Rate}
First, calculate the mass flow rate \(\dot{m}\) using:
\[
\dot{m} = \rho A_1 V_1
\]
Area \(A_1\) at Section (1) is:
\[
A_1 = \pi \left( \frac{D_1}{2} \right)^2 = \pi \left( \frac{0.3 \, \text{m}}{2} \right)^2 = 0.0707 \, \text{m}^2
\]
For water, \(\rho = 1000 \, \text{kg/m}^3\),
\[
\dot{m} = 1000 \, \text{kg/m}^3 \times 0.0707 \, \text{m}^2 \times 2 \, \text{m/s} = 141.4 \, \text{kg/s}
\]

\subsubsection*{Velocity at Section (2)}
To find \(V_2\), we use the equation of continuity \(\dot{m} = \rho A_2 V_2\):
\[
V_2 = \frac{\dot{m}}{\rho A_2}
\]
Area \(A_2\) at Section (2) is:
\[
A_2 = \pi \left( \frac{D_2}{2} \right)^2 = \pi \left( \frac{0.16 \, \text{m}}{2} \right)^2 = 0.0201 \, \text{m}^2
\]
\[
V_2 = \frac{141.4 \, \text{kg/s}}{1000 \, \text{kg/m}^3 \times 0.0201 \, \text{m}^2} = 7.035 \, \text{m/s}
\]

\subsubsection*{Force due to Pressure Change}
The force \( F_{P1} \) due to the pressure at Section (1) is \( P_1 A_1 \) and the force \( F_{P2} \) is 0.
\[
F_{P1} = 100 \times 10^3 \, \text{Pa} \times 0.0707 \, \text{m}^2 = 7070 \, \text{N}
\]
\[
F_{P2} = 0 \, \text{N}
\]
Net force due to pressure change:
\[
\Delta F_{P} = F_{P1} - F_{P2} = 7070 \, \text{N} - 0 \, \text{N} = 7070 \, \text{N}
\]

\subsubsection*{Force due to Velocity Change}
The change in momentum per unit time gives the force \( F_{V} \) as:
\[
F_{V} = \dot{m} (V_2 - V_1) = 141.4 \, \text{kg/s} \times (7.02 \, \text{m/s} - 2 \, \text{m/s}) = 707.14 \, \text{N}
\]
*Note: Chegg says my solution is wrong. The correct solution is: 
\[
F_{V} = \dot{m} (V_2 + V_1) = 141.4 \, \text{kg/s} \times (7.02 \, \text{m/s} + 2 \, \text{m/s}) = 1275.43 \, \text{N}
\]

\subsubsection*{Total Anchoring Force}
Finally, the total anchoring force \( F \) needed is:
\[
F = \Delta F_{P} + F_{V} = 7070 + 707.14 \approx 7777.14 \, \text{N}
\]
*Note: Chegg Solution States: 
\[
F = \Delta F_{P} + F_{V} = 7070 + 1275.43 \approx 8345.43 \, \text{N}
\]





\newpage
\section*{Problem \#3}
Determine the thrust needed from the propeller to hold the hydraulic dredge boat stationary, given:
\begin{itemize}
    \item Depth of intake below surface \(d = 7 \, \text{ft}\)
    \item Total depth of the river \(H = 9 \, \text{ft}\)
    \item Ejection velocity \(V = 30 \, \text{ft/s}\)
    \item Diameter of the ejected stream \(D = 2 \, \text{ft}\)
    \item Angle of ejection \(\theta = 30^\circ\)
    \item Specific gravity \(SG = 1.4\)
    \item Density of water \(\rho_{\text{water}} = 1.94 \, \text{slug/ft}^3\)
\end{itemize}

\subsection*{Assumptions}
\begin{enumerate}
    \item Steady flow
    \item One-dimensional flow
    \item No other forces acting on the fluid element
\end{enumerate}

\subsection*{Calculations}

\subsubsection*{Mass Flow Rate}
First, calculate the mass flow rate \(\dot{m}\):
\[
\dot{m} = \rho A V
\]
Area \(A\) of the ejected stream is:
\[
A = \pi \left( \frac{D}{2} \right)^2 = \pi \left( \frac{2 \, \text{ft}}{2} \right)^2 = 3.142 \, \text{ft}^2
\]
Density of the sand/water mixture \(\rho_{\text{mixture}}\) is:
\[
\rho_{\text{mixture}} = SG \times \rho_{\text{water}} = 1.4 \times 1.94 \, \text{slug/ft}^3 = 2.716 \, \text{slug/ft}^3
\]
\[
\dot{m} = 2.716 \, \text{slug/ft}^3 \times 3.142 \, \text{ft}^2 \times 30 \, \text{ft/s} = 256 \, \text{slug/s}
\]

\subsubsection*{Horizontal Component of Force}
Calculate the horizontal component of the force using \(\dot{m}\) and \(V\):
\[
F_x = \dot{m} \times V \times \cos \theta
\]
\[
F_x = 256 \, \text{slug/s} \times 30 \, \text{ft/s} \times \cos 30^\circ = 256 \, \text{slug/s} \times 30 \, \text{ft/s} \times \frac{\sqrt{3}}{2} \approx 6650 \, \text{lbf}
\]
The thrust needed from the propeller to hold the boat stationary is approximately \(6650 \, \text{lbf}\).





\newpage
\section*{Problem \#4}
Determine the value of the restraining force \( F_x \) exerted on the rocket, given:
\begin{itemize}
    \item Velocity of the exhaust gas \(V = 5000 \, \text{ft/s}\)
    \item Pressure of the exhaust gas \(P = 20 \, \text{psia}\)
    \item Area of the nozzle \(A = 60 \, \text{in}^2\)
    \item Mass flow rate \(\dot{m} = 21 \, \text{lbm/s}\)
\end{itemize}

\subsection*{Assumptions}
\begin{enumerate}
    \item Steady flow
    \item One-dimensional flow
    \item No other forces acting on the exhaust gas
    \item \textbf{We assume there is no atmospheric pressure}
\end{enumerate}

\subsection*{Calculations}
\subsubsection*{Momentum and Pressure Forces}
To find the value of the restraining force \( F_x \), we need to consider both the momentum and the pressure forces in the x-direction:
\[
F_x = \dot{m} \times V + P \times A
\]
The area \( A \) is given in \( \text{in}^2 \), and pressure \( P \) is given in \( \text{psia} \). One \( \text{psia} \) is equal to \( 1 \, \text{lb/in}^2 \), so we can directly use the given units in the formula. Thus, we have:
\[
F_x = 21 \, \text{lbm/s} \times 5000 \, \text{ft/s} + 20 \, \text{psia} \times 60 \, \text{in}^2
\]
\[
F_x = 105000 \, \text{lbm ft/s}^2 + 1200 \, \text{lb}
\]
Now, to convert the momentum force to lb, we use the conversion factor \(1 \, \text{lb} = 32.174 \, \text{lbm ft/s}^2\):
\[
F_x = \frac{105000 \, \text{lbm ft/s}^2}{32.174 \, \text{lbm ft/s}^2/\text{lb}} + 1200 \, \text{lb}
\]
\[
F_x \approx 3263.2 \, \text{lb} + 1200 \, \text{lb} = 4463.2 \, \text{lb}
\]
The value of the restraining force \( F_x \) is approximately \( 4463.2 \, \text{lb} \).
\subsubsection*{Momentum and Pressure Forces Considering Atmospheric Pressure}
To find the value of the restraining force \( F_x \), we consider both the momentum and the net pressure forces in the x-direction:
\[
F_x = \dot{m} \times V + (P - P_{\text{atm}}) \times A
\]
Where \( P_{\text{atm}} \) is the atmospheric pressure, \( 14.7 \, \text{psia} \). Thus, the formula becomes:
\[
F_x = 21 \, \text{lbm/s} \times 5000 \, \text{ft/s} + (20 \, \text{psia} - 14.7 \, \text{psia}) \times 60 \, \text{in}^2
\]

\[
F_x = \frac{105000 \, \text{lbm ft/s}^2}{32.174 \, \text{lbm ft/s}^2/\text{lb}} + 318 \, \text{lb}
\]
\[
F_x \approx 3263.2 \, \text{lb} + 318 \, \text{lb} = 3581.2 \, \text{lb}
\]
The value of the restraining force \( F_x \) is approximately \( 3581.2 \, \text{lb} \).




\newpage
\section*{Problem \#5}
Determine the density of the exhaust gases, given:
\begin{itemize}
    \item Air mass flow rate \(\dot{m}_{\text{air}} = 65 \, \text{lbm/s}\)
    \item Fuel mass flow rate \(\dot{m}_{\text{fuel}} = 0.60 \, \text{lbm/s}\)
    \item Exhaust gas velocity \(V = 1500 \, \text{ft/s}\)
    \item Effective cross-sectional area \(A = 3.5 \, \text{ft}^2\)
\end{itemize}

\subsection*{Assumptions}
\begin{enumerate}
    \item Steady flow
    \item One-dimensional flow
    \item No other forces acting on the fluid element
\end{enumerate}

\subsection*{Calculations}

\subsubsection*{Total Mass Flow Rate}
First, calculate the total mass flow rate \(\dot{m}_{\text{total}}\):
\[
\dot{m}_{\text{total}} = \dot{m}_{\text{air}} + \dot{m}_{\text{fuel}} = 65 \, \text{lbm/s} + 0.60 \, \text{lbm/s} = 65.60 \, \text{lbm/s}
\]

\subsubsection*{Density of the Exhaust Gases}
Use the formula \(\rho = \frac{\dot{m}}{A \times V}\) to find the density of the exhaust gases:
\[
\rho = \frac{65.60 \, \text{lbm/s}}{3.5 \, \text{ft}^2 \times 1500 \, \text{ft/s}} = \frac{65.60 \, \text{lbm/s}}{5250 \, \text{ft}^3/\text{s}} = 0.0125 \, \text{lbm/ft}^3
\]
The estimated density of the exhaust gases is \(0.0125 \, \text{lbm/ft}^3\).




\end{document}
